% Distilled blueprint for Evans, Appendix E (Measure Theory), including the
% concluding spectral theorem from Appendix D that occurs on the opening page.

\begin{theorem}[Eigenvectors of a compact, symmetric operator]
\label{thm:eigenvectors-compact-symmetric-operator}
\uses{def:hilbert-space, def:separable-space, def:compact-linear-operator, def:adjoint-symmetric-operator, def:eigenvalue-eigenvector-appendix-d, thm:fredholm-alternative, thm:spectrum-of-compact-operator, lem:bounds-on-spectrum-symmetric-operator}
\dcref{appD:thm-7}
\group{Fredholm Theory}

Let $H$ be a separable Hilbert space. Every compact symmetric operator
$S:H\to H$ has a countable orthonormal basis of eigenvectors.
\end{theorem}
\begin{proof}
\sketch List the distinct nonzero eigenvalues of $S$ as $\{\eta_k\}_{k\geq1}$,
set $\eta_0=0$, and write
\[
H_0=N(S),\qquad H_k=N(S-\eta_kI)\quad(k\geq1).
\]
The Fredholm alternative makes every $H_k$, $k\geq1$, nonzero and finite
dimensional; $H_0$ may have any dimension. If $u\in H_k$, $v\in H_l$, and
$k\ne l$, symmetry gives
\[
\eta_k(u,v)=(Su,v)=(u,Sv)=\eta_l(u,v),
\]
so the eigenspaces are mutually orthogonal.

Let $V$ be their algebraic span. Both $V$ and $V^\perp$ are invariant under
$S$. The restriction $\widetilde S=S|_{V^\perp}$ is compact and symmetric and
has no nonzero eigenvalue; the compact-spectrum theorem therefore gives it no
nonzero spectral value. The bounds on the spectrum of a symmetric operator
imply $(\widetilde Su,u)=0$ for $u\in V^\perp$. Polarization yields
\[
2(\widetilde Su,v)
=(\widetilde S(u+v),u+v)-(\widetilde Su,u)-(\widetilde Sv,v)=0,
\]
hence $\widetilde S=0$. Thus $V^\perp\subset N(S)=H_0\subset V$, so
$V^\perp=\{0\}$ and $V$ is dense in $H$. Choose an orthonormal basis of each
$H_k$; separability makes the basis of $H_0$ countable, and their union is the
required basis of $H$.
\end{proof}

Proof references include \cite{Brezis1983}, \cite[Chapter 5]{GilbargTrudinger1983},
and \cite{Yosida1980}.

\section{Measure Theory}
\label{sec:measure-theory}

\subsection{Lebesgue Measure}
\label{sec:lebesgue-measure}

\begin{definition}[$\sigma$-algebra]
\label{def:sigma-algebra}
\dcref{appE:E.1-def-1}
\group{Measure Theory}
A collection $\mathcal M$ of subsets of $\mathbb R^n$ is a $\sigma$-algebra if:

(i) $\emptyset,\mathbb R^n\in\mathcal M$;

(ii) $A\in\mathcal M$ implies $\mathbb R^n-A\in\mathcal M$;

(iii) whenever $\{A_k\}_{k=1}^\infty\subset\mathcal M$, both
$\bigcup_{k=1}^\infty A_k$ and $\bigcap_{k=1}^\infty A_k$ belong to
$\mathcal M$.
\end{definition}

\begin{theorem}[Existence of Lebesgue measure and Lebesgue measurable sets]
\label{thm:existence-lebesgue-measure}
\uses{def:sigma-algebra}
\dcref{appE:thm-1}
\group{Measure Theory}

There are a $\sigma$-algebra $\mathcal M$ of subsets of $\mathbb R^n$ and a map
\[
|\cdot|:\mathcal M\to[0,+\infty]
\]
such that:

(i) every open, and hence every closed, subset of $\mathbb R^n$ lies in
$\mathcal M$;

(ii) for every ball $B\subset\mathbb R^n$, $|B|$ is its $n$-dimensional
volume;

(iii) if $\{A_k\}_{k=1}^\infty\subset\mathcal M$ is pairwise disjoint, then
\[
(1)\qquad
\left|\bigcup_{k=1}^\infty A_k\right|=\sum_{k=1}^\infty|A_k|;
\]

(iv) if $A\subseteq B$, $B\in\mathcal M$, and $|B|=0$, then
$A\in\mathcal M$ and $|A|=0$.

The members of $\mathcal M$ are the Lebesgue measurable sets, and $|\cdot|$ is
$n$-dimensional Lebesgue measure.
\end{theorem}

\begin{remark}[Basic properties of Lebesgue measure]
\label{rem:lebesgue-measure-properties}
\uses{thm:existence-lebesgue-measure}
\dcref{appE:E.1-rem-1}
\group{Measure Theory}

(i) The measure of a set with piecewise smooth boundary equals its volume.

(ii) Countable additivity implies
\[
(2)\qquad |\emptyset|=0
\]
and, for every countable family of measurable sets,
\[
(3)\qquad
\left|\bigcup_{k=1}^\infty A_k\right|\leq\sum_{k=1}^\infty|A_k|.
\]
\end{remark}

\begin{definition}[Almost everywhere]
\label{def:almost-everywhere}
\uses{thm:existence-lebesgue-measure}
\dcref{appE:E.1-def-2}
\group{Measure Theory}

A property holds \textit{almost everywhere}, abbreviated a.e., if it fails only
on a measurable set of Lebesgue measure zero.
\end{definition}

\subsection{Measurable Functions and Integration}
\label{sec:measurable-functions-integration}

\begin{definition}[Measurable function]
\label{def:measurable-function}
\uses{def:sigma-algebra}
\dcref{appE:E.2-def-1}
\group{Measure Theory}

A function $f:\mathbb R^n\to\mathbb R$ is measurable if
\[
f^{-1}(U)\in\mathcal M
\]
for every open set $U\subset\mathbb R$.
\end{definition}

Continuous functions are measurable. Sums and products of measurable functions
are measurable, as are the pointwise $\limsup$ and $\liminf$ of any sequence of
measurable functions.

\begin{theorem}[Egoroff's Theorem]
\label{thm:egoroffs-theorem}
\uses{def:measurable-function, def:almost-everywhere}
\dcref{appE:thm-2}
\group{Measure Theory}

Let $A\subset\mathbb R^n$ be measurable with $|A|<\infty$, and let the
measurable functions $f_k,f$ satisfy $f_k\to f$ a.e. on $A$. For every
$\varepsilon>0$ there is a measurable set $E\subset A$ such that
\[
|A-E|\leq\varepsilon
\]
and $f_k\to f$ uniformly on $E$.
\end{theorem}

For a nonnegative measurable function $f$, its Lebesgue integral
\[
\int_{\mathbb R^n}f\,dx
\]
is defined by approximation with simple functions; compare
$\S\ref{sec:banach-space-valued-functions}$. It agrees with the usual integral
for continuous or Riemann-integrable functions. For an arbitrary measurable
function, define
\[
\int_{\mathbb R^n}f\,dx
:=\int_{\mathbb R^n}f^+\,dx-\int_{\mathbb R^n}f^-\,dx
\]
provided at least one term on the right is finite. Such a function is called
integrable.

\begin{definition}[Summable function]
\label{def:summable-function}
\uses{def:measurable-function}
\dcref{appE:E.2-def-2}
\group{Measure Theory}

A measurable function $f$ is summable if
\[
\int_{\mathbb R^n}|f|\,dx<\infty.
\]
\end{definition}

\begin{remark}[Integrable versus summable terminology]
\label{rem:integrable-vs-summable-terminology}
\uses{def:summable-function}
\dcref{appE:E.2-rem-1}
\group{Measure Theory}

Here integrable means that the integral is defined, possibly with value
$+\infty$ or $-\infty$, whereas summable means that the integral is finite.
\end{remark}

\begin{definition}[Essential supremum]
\label{def:essential-supremum}
\dcref{appE:E.2-def-3}
\group{Measure Theory}
For a measurable real-valued function $f$,
\[
\operatorname*{ess\,sup}f
:=\inf\{\mu\in\mathbb R\mid |\{f>\mu\}|=0\}.
\]
\end{definition}

\subsection{Convergence Theorems for Integrals}
\label{sec:convergence-theorems-for-integrals}

\begin{theorem}[Fatou's Lemma]
\label{thm:fatous-lemma}
\uses{def:summable-function}
\dcref{appE:thm-3}
\group{Measure Theory}

Let $\{f_k\}_{k=1}^\infty$ and $f$ be nonnegative summable functions. Then
\[
\int_{\mathbb R^n}\liminf_{k\to\infty}f_k\,dx
\leq\liminf_{k\to\infty}\int_{\mathbb R^n}f_k\,dx.
\]
\end{theorem}

\begin{theorem}[Monotone Convergence Theorem]
\label{thm:monotone-convergence-theorem}
\uses{def:summable-function}
\dcref{appE:thm-4}
\group{Measure Theory}

Let $\{f_k\}_{k=1}^\infty$ be measurable and satisfy
\[
f_1\leq f_2\leq\cdots\leq f_k\leq f_{k+1}\leq\cdots.
\]
Then
\[
\int_{\mathbb R^n}\lim_{k\to\infty}f_k\,dx
=\lim_{k\to\infty}\int_{\mathbb R^n}f_k\,dx.
\]
\end{theorem}

\begin{theorem}[Dominated Convergence Theorem]
\label{thm:dominated-convergence-theorem}
\uses{def:summable-function}
\dcref{appE:thm-5}
\group{Measure Theory}

Let $\{f_k\}_{k=1}^\infty$ be integrable, suppose $f_k\to f$ a.e., and assume
that a summable function $g$ satisfies
\[
|f_k|\leq g\quad\text{a.e.}
\]
for every $k$. Then
\[
\int_{\mathbb R^n}f_k\,dx\to\int_{\mathbb R^n}f\,dx.
\]
\end{theorem}

\subsection{Differentiation}
\label{sec:differentiation}

\begin{theorem}[Lebesgue's Differentiation Theorem]
\label{thm:lebesgue-differentiation-theorem}
\uses{def:summable-function, def:almost-everywhere}
\dcref{appE:thm-6}
\group{Measure Theory}

Let $f:\mathbb R^n\to\mathbb R$ be locally summable.

(i) For a.e. $x_0\in\mathbb R^n$,
\[
\fint_{B(x_0,r)}f\,dx\to f(x_0)\qquad(r\to0).
\]

(ii) More precisely, for a.e. $x_0\in\mathbb R^n$,
\[
(4)\qquad
\fint_{B(x_0,r)}|f(x)-f(x_0)|\,dx\to0
\qquad(r\to0).
\]

A point $x_0$ satisfying (4) is a Lebesgue point of $f$.
\end{theorem}

\begin{remark}[$L^p$ version of the Lebesgue point theorem]
\label{rem:lebesgue-differentiation-lp}
\uses{thm:lebesgue-differentiation-theorem}
\dcref{appE:E.4-rem-1}
\group{Measure Theory}

If $f\in L^p_{\mathrm{loc}}(\mathbb R^n)$ for $1\leq p<\infty$, then for a.e.
$x_0\in\mathbb R^n$,
\[
\fint_{B(x_0,r)}|f(x)-f(x_0)|^p\,dx\to0
\qquad(r\to0).
\]
\end{remark}

\subsection{Banach Space-Valued Functions}
\label{sec:banach-space-valued-functions}

Let $T>0$, let $X$ be a real Banach space with norm $\|\cdot\|$, and consider
functions $f:[0,T]\to X$.

\begin{definition}[Simple function]
\label{def:simple-function-banach-valued}
\uses{def:sigma-algebra}
\dcref{appE:E.5-def-1}
\group{Vector-Valued Integration}

A function $s:[0,T]\to X$ is simple if
\[
(5)\qquad
s(t)=\sum_{i=1}^m\chi_{E_i}(t)u_i
\qquad(0\leq t\leq T),
\]
where each $E_i\subset[0,T]$ is Lebesgue measurable and $u_i\in X$.
\end{definition}

\begin{definition}[Strongly measurable function]
\label{def:strongly-measurable-function}
\uses{def:simple-function-banach-valued, def:almost-everywhere}
\dcref{appE:E.5-def-2}
\group{Vector-Valued Integration}

A function $f:[0,T]\to X$ is strongly measurable if simple functions
$s_k:[0,T]\to X$ exist such that
\[
s_k(t)\to f(t)\quad\text{for a.e. }t\in[0,T].
\]
\end{definition}

\begin{definition}[Weakly measurable function]
\label{def:weakly-measurable-function}
\uses{def:measurable-function, def:bounded-linear-functional-dual-space}
\dcref{appE:E.5-def-3}
\group{Vector-Valued Integration}

A function $f:[0,T]\to X$ is weakly measurable if
$t\mapsto\langle u^*,f(t)\rangle$ is Lebesgue measurable for every $u^*\in X^*$.
\end{definition}

\begin{definition}[Almost separably valued]
\label{def:almost-separably-valued}
\uses{def:separable-space}
\dcref{appE:E.5-def-4}
\group{Vector-Valued Integration}

A function $f:[0,T]\to X$ is almost separably valued if there is a set
$N\subset[0,T]$ with $|N|=0$ such that
\[
\{f(t)\mid t\in[0,T]-N\}
\]
is separable.
\end{definition}

\begin{theorem}[Pettis]
\label{thm:pettis-theorem}
\uses{def:strongly-measurable-function, def:weakly-measurable-function, def:almost-separably-valued}
\dcref{appE:thm-7}
\group{Vector-Valued Integration}

A map $f:[0,T]\to X$ is strongly measurable if and only if it is weakly
measurable and almost separably valued.
\end{theorem}

\begin{definition}[Integral of a simple function]
\label{def:integral-of-simple-function}
\uses{def:simple-function-banach-valued}
\dcref{appE:E.5-def-5}
\group{Vector-Valued Integration}

For a simple function $s(t)=\sum_{i=1}^m\chi_{E_i}(t)u_i$, define
\[
(6)\qquad
\int_0^T s(t)\,dt:=\sum_{i=1}^m|E_i|u_i.
\]
\end{definition}

\begin{definition}[Summable Banach space-valued function]
\label{def:summable-banach-valued-function}
\uses{def:integral-of-simple-function}
\dcref{appE:E.5-def-6}
\group{Vector-Valued Integration}

A function $f:[0,T]\to X$ is summable if there are simple functions
$\{s_k\}_{k=1}^\infty$ such that
\[
(7)\qquad
\int_0^T\|s_k(t)-f(t)\|\,dt\to0
\qquad(k\to\infty).
\]
\end{definition}

\begin{definition}[Bochner integral]
\label{def:bochner-integral}
\uses{def:summable-banach-valued-function, def:integral-of-simple-function}
\dcref{appE:E.5-def-7}
\group{Vector-Valued Integration}

For a summable $f$ and any sequence of simple functions satisfying (7), define
\[
(8)\qquad
\int_0^T f(t)\,dt:=\lim_{k\to\infty}\int_0^T s_k(t)\,dt.
\]
\end{definition}

\begin{theorem}[Bochner]
\label{thm:bochner-theorem}
\uses{def:strongly-measurable-function, def:summable-banach-valued-function, def:bochner-integral, def:summable-function}
\dcref{appE:thm-8}
\group{Vector-Valued Integration}

A strongly measurable function $f:[0,T]\to X$ is summable if and only if the
scalar function $t\mapsto\|f(t)\|$ is summable. In that case,
\[
\left\|\int_0^T f(t)\,dt\right\|
\leq\int_0^T\|f(t)\|\,dt
\]
and, for every $u^*\in X^*$,
\[
\left\langle u^*,\int_0^T f(t)\,dt\right\rangle
=\int_0^T\langle u^*,f(t)\rangle\,dt.
\]
\end{theorem}

A general measure-theory reference is \cite{Folland1984}. Proofs of
\cref{thm:pettis-theorem} and \cref{thm:bochner-theorem} are given in
\cite[Chapter V, Sections 4--5]{Yosida1980}.
